\subsection{Value functions inherit curvature and growth}\label{sec:valuefn}
Split the coordinates into \emph{retained} coordinates $v=x_K$ and
\emph{recourse} coordinates $y=x_R$, $R=[n]\setminus K$, and define
\begin{equation}\label{eq:valuefn}
V(v)=\min_{y\in X_R}F(v,y),\qquad v\in X_K .
\end{equation}
The feasible set $X_R$ of the recourse problem does not depend on $v$; this
is the only structural requirement. Clearly $\min_{X_K}V=\OPT$, and a
retained point $v$ together with a minimizer $y(v)$ of the recourse problem
is a feasible point of~\eqref{eq:model} with value $V(v)$.

\begin{lemma}[Partial minimization]\label{lem:valuefunction}
Let $V$ be defined by~\eqref{eq:valuefn}, with $X_R$ compact.
\begin{enumerate}[label=(\alph*)]
\item If $L_i$ is an upper coordinate curvature of $F$ for $i\in K$, it is
an upper coordinate curvature of $V$ on the continuous hull of $X_K$.
\item If $F$ has quadratic growth with constant $g$ at $(v^*,y^*)$, then $V$
has quadratic growth with constant $g$ at $v^*$ on $X_K$.
\end{enumerate}
\end{lemma}

\begin{proof}
(a) For fixed $v$, $i\in K$ and $y\in X_R$, the function
$t\mapsto F(v+te_i,y)-\frac{L_i}{2}t^2$ is concave, and
$t\mapsto V(v+te_i)-\frac{L_i}{2}t^2$ is the pointwise minimum of these
functions over $y\in X_R$; a pointwise infimum of concave functions is
concave. (b) For $v\in X_K$ let $y(v)$ attain the minimum. Then
$V(v)-V(v^*)=F(v,y(v))-\OPT\ge g(\norm{v-v^*}^2+\norm{y(v)-y^*}^2)\ge
g\norm{v-v^*}^2$.
\end{proof}

For a nonempty set $B$ of coordinates we write $W_B$ for the value function
\eqref{eq:valuefn} with $K=B$, and $\bar X_B=\prod_{i\in B}[\ell_i,u_i]$.
A \emph{bag cell} is a product $C=\prod_{i\in B}[a_i,b_i]$ with
$a_i\le b_i$ and $a_i,b_i\in X_i$. Its \emph{corners} are the points $v$
with $v_i\in\{a_i,b_i\}$; they are feasible bag values. Put
\[
e_B(C)=\frac L8\sum_{i\in B}w_i(C)^2,\qquad
w_i(C)=\begin{cases}
0,& i\in I_Z\text{ and }b_i-a_i=1,\\
b_i-a_i,&\text{otherwise.}
\end{cases}
\]
If every side of $C$ is at most $h$, then $e_B(C)\le|B|Lh^2/8$.

\begin{lemma}[Conditional values inherit coordinate curvature]\label{lem:cr-semiconcave}
For every $i\in B$ and every fixed $v_{B\setminus\{i\}}$, the function
$t\mapsto W_B(v_{B\setminus\{i\}},t)-Lt^2/2$ is concave on $[l_i,u_i]$.
\end{lemma}
\begin{proof}
Fix $x_{-B}\in X_{-B}$. The point $(v_{B\setminus\{i\}},t,x_{-B})$ lies in
the continuous hull of $X$ for $t\in[l_i,u_i]$, so
$t\mapsto F(v_{B\setminus\{i\}},t,x_{-B})-Lt^2/2$ is concave. The
function in the statement is the pointwise minimum of these concave
functions over the set $X_{-B}$, which does not depend on $t$. A pointwise
infimum of concave functions is concave, because its hypograph is the
intersection of their hypographs, and here it is finite.
\end{proof}

\begin{lemma}[Bag-cell interpolation]\label{lem:cr-cell}
Let $C$ be a bag cell. Every $x\in X$ with $x_B\in C$ satisfies
\[
F(x)\ \ge\ \min_{v\text{ a corner of }C}W_B(v)-e_B(C).
\]
\end{lemma}
\begin{proof}
Put $\xi=x_B$. Let $(Y_i)_{i\in B}$ be independent, with $Y_i\in\{a_i,b_i\}$
and $\E Y_i=\xi_i$ (and $Y_i=a_i$ if $a_i=b_i$). Then
$\operatorname{Var}Y_i=(\xi_i-a_i)(b_i-\xi_i)\le(b_i-a_i)^2/4$. If $i$ is an
integer coordinate with $b_i-a_i=1$, then $\xi_i\in\{a_i,b_i\}$ because
$x$ is feasible, and $\operatorname{Var}Y_i=0$. Hence
$\frac L2\sum_i\operatorname{Var}Y_i\le e_B(C)$.

Order $B$ as $i_1,\ldots,i_r$, and let $Z^{(s)}$ be $\xi$ with its first $s$
listed coordinates replaced by the corresponding $Y$'s. All $Z^{(s)}$ lie
in $\bar X_B$. Conditionally on $Y_{i_1},\ldots,Y_{i_{s-1}}$, the
function $\varphi(t)=W_B(Z^{(s-1)}\text{ with coordinate }i_s\text{ set to }t)$
has $\varphi(t)-Lt^2/2$ concave by Lemma~\ref{lem:valuefunction}(a). Jensen's
inequality for this concave function gives
$\E\varphi(Y_{i_s})\le\varphi(\xi_{i_s})+\frac L2\operatorname{Var}Y_{i_s}$,
that is,
$\E[W_B(Z^{(s)})\mid Y_{i_1},\ldots,Y_{i_{s-1}}]
\le W_B(Z^{(s-1)})+\frac L2\operatorname{Var}Y_{i_s}$.
Taking expectations and summing over $s$ yields
$\E W_B(Y)\le W_B(\xi)+e_B(C)$. Since $Y$ is a corner of $C$ almost surely,
\[
\min_{\text{corners}}W_B\le\E W_B(Y)\le W_B(\xi)+e_B(C)\le F(x)+e_B(C).
\qedhere
\]
\end{proof}

No coordinate outside $B$ is rounded, so the error $e_B(C)$ does not grow
with $n$. The price is that $W_B$ is a global optimization problem over
the original outside domain.


Part~(a) explains why the corrected-grid bound survives partial
minimization, and part~(b) why the growth analysis does. The curvature of
$V$ can be much smaller than that of $F$. If $F(v,y)=M(y-v)^2$ with
$v,y\in[0,1]$, the diagonal curvature in $v$ is $2M$, whereas $V\equiv0$.
Sections~\ref{sec:convexrecourse} and~\ref{sec:cuts} give certified ways to
evaluate $V$ and to bound its curvature by quantities that do not see such
stiff convex terms.

The decomposition of~\eqref{eq:model} is used as follows. Suppose the
recourse coordinates form disjoint \emph{private blocks} $y^{(1)},\dots,
y^{(r)}$ and
\begin{equation}\label{eq:privateblocks}
F(v,y)=F_0(v)+\sum_{t=1}^r\phi_t\bigl(v_{E_t},y^{(t)}\bigr),
\end{equation}
where $F_0$ is a sum of factors in $v$ and each block interacts with the
retained coordinates only through the scope $E_t\subseteq K$. Then
\[
V(v)=F_0(v)+\sum_{t=1}^r\psi_t(v_{E_t}),\qquad
\psi_t(w)=\min_{y^{(t)}}\phi_t(w,y^{(t)}),
\]
is again a sum of factors, the \emph{value factors} $\psi_t$, with scopes
$E_t$. Any tree decomposition of the scopes of $F_0$ and the sets $E_t$
serves for $V$. Its width is the width after elimination of the blocks,
which can differ from the width of the original decomposition.

\begin{theorem}[Grids over value functions]\label{thm:valuefn}
Suppose $V$ has the form above with a tree decomposition of maximum bag
size $p$, rational upper coordinate curvatures $L_i^V$ ($i\in K$), and
value factors that can be evaluated exactly at rational points in time
polynomial in the input and the bit length of the point. Then
Theorem~\ref{thm:certificate} holds for $V$, and if $F$ has quadratic
growth, Theorem~\ref{thm:approx} holds with $\kappa_V=\max\{1,L^V/g\}$ in
place of $\kappa$, where $L^V=\max_iL_i^V$. If $F$ is a rational quadratic
and each recourse problem can be solved exactly, Theorem~\ref{thm:exact}
holds with $\kappa_V$ as well.
\end{theorem}

\begin{proof}
Sections~\ref{sec:grids} and~\ref{sec:growth} use the objective only
through Definition~\ref{def:curvature}, exact evaluation of factors at
grid points, and growth; by Lemma~\ref{lem:valuefunction} these hold for
$V$. In the bit analysis of Theorem~\ref{thm:approx}, a table entry is a
sum of at most $\abs{\mathcal A}+n$ exactly evaluated factor values, each of
polynomial bit length, so the arithmetic bounds are unchanged. For exact
output, the incumbent is a feasible point $(\hat v,y(\hat v))$ of the
original problem whose value is within the certified gap of $\OPT$, and
growth of $F$ is set growth with $S=\{(v^*,y^*)\}$; Theorem~\ref{thm:transfer}
applied to the original quadratic therefore turns it into an exact
minimizer once the gap is below $\varepsilon_S$, which happens after
$\poly(I)+O(\log\kappa_V)$ bits of accuracy because $g\ge L^V/\kappa_V$.
\end{proof}

The theorem charges the work of the recourse oracle to the polynomial
factor, so it is useful when the recourse problems are easy and their
elimination does not increase the width. The next two subsections give
classes where both hold and the curvature bound $L^V$ is certified.
